% Generated draft: requires mathematical and source review.
\documentclass[11pt]{article}
\usepackage[margin=25mm]{geometry}
\usepackage{fontspec}
\setmainfont{Latin Modern Roman}
\usepackage{amsmath,amssymb,mathtools,unicode-math}
\setmathfont{Latin Modern Math}
\usepackage{enumitem,needspace,xurl,hyperref}
\setlength{\parindent}{0pt}
\setlength{\parskip}{5pt}
\setlength{\emergencystretch}{3em}
\newcommand{\N}{\mathbb{N}}
\newcommand{\Z}{\mathbb{Z}}
\newcommand{\Q}{\mathbb{Q}}
\newcommand{\R}{\mathbb{R}}
\newcommand{\C}{\mathbb{C}}
\newcommand{\F}{\mathbb{F}}
\newcommand{\A}{\mathbb{A}}
\newcommand{\PP}{\mathbb{P}}
\newcommand{\E}{\mathbb{E}}
\newcommand{\eps}{\varepsilon}
\newcommand{\dd}{\,\mathrm{d}}
\begin{document}
\section*{UnsolvedMath ID: 11000110}
\textbf{Source problem code:} catalogue entry
\subsection{Definitions and mathematical statement}
Let $g>h$ with $g\geq3$, let $\Gamma\leq\operatorname{Mod}_g$ be any finite-index subgroup, and let $\varphi:\Gamma\to\operatorname{Mod}_h$ be any homomorphism. Problem 2.5 asks:
\begin{enumerate}
\item[(a)] Must $\varphi(\Gamma)$ be finite?
\item[(b)] Must $[\operatorname{Mod}_h:\varphi(\Gamma)]$ be infinite?
\end{enumerate}
Both questions quantify over all such subgroups and homomorphisms, with no normality or faithfulness requirement. The source calls (b) weaker than (a); that implication requires the target mapping class group to be infinite. In the closed unmarked convention, the genus-zero target is trivial, so its inclusion would make (b) false. The omitted genus and surface conventions must therefore be checked; the written assumptions do not explicitly exclude $h=0$.

The following extracted paragraph turns to a related virtual first-cohomology problem. Its explanation identifies the invariant as
\[
 H^1(\operatorname{Mod}_g;\Z)
 =\operatorname{Hom}(\operatorname{Mod}_g,\Z),
\]
reported to vanish, and asks whether the same holds for finite-index subgroups, attributing this to Ivanov [25], Problem 2.11(A). First cohomology, rather than vanishing integral abelianization $H_1$, is the correct interpretation of the displayed explanation; those assertions differ. This context is not a replacement for either of the two homomorphism questions above.
\subsection{Short English statement}
Must every homomorphism from a finite-index genus-$g$ mapping-class subgroup to a lower-genus mapping class group have finite image, or at least infinite-index image? The target's low-genus convention is essential.
\subsection{Sources}
\begin{itemize}
\item Local registry: \url{data/unsolvedmath/problems.json}, permanent UnsolvedMath ID $11000110$, exact source code \texttt{catalogue entry}.
\item Farb, \emph{Problems on Mapping Class Groups and Related Topics}, Problem 2.5, PDF page 94: \url{https://www.math.uchicago.edu/~farb/papers/mcgbook.pdf}. Supplied access date: 2026-07-29; extraction: pdf-font-structure; repeated local numbering disambiguated by page among 329 headings.
\item The following paragraph cites Ivanov [25], Problem 2.11(A), for the finite-index first-cohomology question. The bibliography entry and ensuing numbered statement are absent. The registry's literature summary is truncated.
\end{itemize}
\end{document}
